% Satellite en orbite circulaire : force, vitesse et accélération dans le repère de Frenet (L'essentiel, chapitre 16)
\begin{tikzpicture}[x=1cm,y=1cm]
\shade[ball color=otNavy!35] (0,0) circle (0.9);
\node[pcLegende,text=white] at (0,0) {Terre};
\draw[otMuted,line width=0.6pt,dash pattern=on 3pt off 2pt] (0,0) circle (2.6);
\coordinate (S) at (35:2.6);
\fill[otInk] (S) circle (0.08);
\node[pcLegende,anchor=south west] at (S) {satellite};
\draw[pcTraj] (0,0) -- (S) node[pos=0.62,pcLegende,anchor=south east] {$r$};
\draw[pcVitesse] (S) -- ++(125:1.6) node[pcLegende,anchor=south,text=otNavy] {$\vec{v}$};
\draw[pcForce] (S) -- ++(215:1.15) node[pcLegende,anchor=north,text=otGreenDark,yshift=-1pt] {$\vec{F}=m\,\vec{a}$};
\draw[otMuted,line width=0.4pt,-{Stealth[length=1.5mm]}] (-60:3.0) arc[start angle=-60,end angle=-20,radius=3.0];
\node[pcLegende,text=otMuted,anchor=north west] at (-40:3.05) {sens de rotation};
\end{tikzpicture}
\par\smallskip\legende{La force de gravitation, dirigée vers le centre de la Terre, est perpendiculaire à la vitesse : elle ne change que la direction du vecteur vitesse.}
